{\footnotesize
\setlength{\tabcolsep}{3pt}\renewcommand{\arraystretch}{1.22}
\begin{longtable}{@{}p{0.42cm}p{4.55cm}p{2.15cm}p{2.45cm}>{\raggedleft\arraybackslash}p{1.45cm}>{\raggedleft\arraybackslash}p{2.0cm}@{}}
\toprule
\rowcolor{rpfnavy!12}
\textbf{ที่} & \textbf{ชื่อโครงการ} & \textbf{ผู้รับผิดชอบ} & \textbf{หน่วยงาน} & \textbf{จัดสรร} & \textbf{เบิกจ่าย}\\
\midrule
\endfirsthead
\toprule
\rowcolor{rpfnavy!12}
\textbf{ที่} & \textbf{ชื่อโครงการ} & \textbf{ผู้รับผิดชอบ} & \textbf{หน่วยงาน} & \textbf{จัดสรร} & \textbf{เบิกจ่าย}\\
\midrule
\endhead
\multicolumn{6}{@{}l}{\cellcolor{rpfnavy!12}\bfseries ง8-1 ผลักดันเทคโนโลยี นวัตกรรมสู่ชุมชน ตามเป้าหมายการพัฒนาอย่างยั่งยืน (16 โครงการ)}\\
\midrule
1 & โครงการการวิจัย Plant profile ของผักในโรงเรือนอัจฉริยะ & นายวัชระ กิตติวรเชฏฐ์ & กลุ่มแผนงานใต้ร่มพระบารมี & 200,000 & 66,387 (33\%)\\
2 & โครงการการวิจัยพัฒนาวิธีการจัดการกระบวนการผลิตผักเพื่อคุณภาพ & นายวัชระ กิตติวรเชฏฐ์ & กลุ่มแผนงานใต้ร่มพระบารมี & 150,000 & 61,222 (41\%)\\
3 & โครงการพัฒนากระบวนการสกัดน้ำมันอะโวคาโดของศูนย์พัฒนาโครงการหลวงห้วยเสี & นางสาวพิมลพรรณ เลิศบัวบา & กลุ่มแผนงานใต้ร่มพระบารมี & 148,500 & 63,002 (42\%)\\
4 & โครงการยกระดับการพัฒนาเทคโนโลยี เครื่องมือ และเครื่องจักรของมหาวิทยาลั & นายสามารถ สาลี & กลุ่มแผนงานใต้ร่มพระบารมี & 52,000 & 0 (0\%)\\
5 & โครงการการพัฒนาตู้อบลมร้อนดอกเก๊กฮวยระบบปั๊มความร้อน เพื่อเพิ่มคุณภาพผ & ผศ.ไกรลาศ ดอนชัย & คณะวิศวกรรมศาสตร์ & 175,000 & 87,500 (50\%)\\
6 & โครงการพัฒนาเทคโนโลยีเครื่องล้างเมล็ดงาดำแบบกึ่งอัตโนมัติ สำหรับศูนย์พ & ผศ.ขัติพงษ์ จิโนสุวัตร์ & คณะวิศวกรรมศาสตร์ & 150,000 & 67,432 (45\%)\\
7 & โครงการพัฒนาเทคโนโลยีต้นแบบเครื่องนวดงาขี้ม่อน(งาหอม) เพื่อเพิ่มประสิท & อ.วรเชษฐ์ หวานเสียง & คณะวิศวกรรมศาสตร์ & 175,000 & 162,092 (93\%)\\
8 & โครงการประเมินสมบัติเชิงกลของเส้นใยกัญชงจากการอบแห้งด้วยเทคโนโลยีอบลมร & นายธีระยุทธ ขอดแก้ว & คณะวิศวกรรมศาสตร์ & 50,000 & 23,200 (46\%)\\
9 & โครงการศึกษาความเหมาะสมของการนำเศษผักเหลือใช้มาผลิตเป็นถ่านอัดแท่งขนาด & นายเมธัส ภัททิยธนี & คณะวิศวกรรมศาสตร์ & 88,000 & 0 (0\%)\\
10 & โครงการ Royal Rose Sensation Giftset ชุดผลิตภัณฑ์ต่อยอดคุณค่ากุหลาบจาก & อ.ณัฐธินี สาลี & วิทยาลัยเทคโนโลยีและสหวิทยาการ & 200,000 & 185,349 (93\%)\\
11 & โครงการพัฒนากระบวนการผลิตและเครื่องมือการอบแห้งเคพกูสเบอร์รี เพื่อเพิ่ & นายวธัญญู วรรณพรหม & สถาบันถ่ายทอดเทคโนโลยีสู่ชุมชน & 100,000 & 20,000 (20\%)\\
12 & โครงการพัฒนาหลักสูตรชุมชนการเรียนรู้เพื่อการพัฒนาตนเองตามภูมิสั่งคมอย่ & ผู้ช่วยศาสตราจารย์ไพโรจน & สถาบันถ่ายทอดเทคโนโลยีสู่ชุมชน & 68,000 & 32,200 (47\%)\\
13 & โครงการพัฒนาชุดทดสอบสารพาราควอตเบื้องต้น (Paraquat Test Kit) ที่เหมาะส & ผศ.สุบิน ใจทา & มทร.ล้านนา เชียงราย & 150,000 & 110,000 (73\%)\\
14 & โครงการนวัตกรรมเชิงพื้นที่ของกระถางเพาะเมล็ดชีวภาพย่อยสลายได้จากของเหล & ผศ.อมรรัตน์ ปิ่นชัยมูล & มทร.ล้านนา เชียงราย & 100,000 & 60,268 (60\%)\\
15 & โครงการประยุกต์ใช้เทคโนโลยีที่เหมาะสมเพื่อเพิ่มประสิทธิภาพการเลี้ยงปลา & นางสาวอพิศรา หงส์หิรัญ & มทร.ล้านนา พิษณุโลก & 100,000 & 77,684 (78\%)\\
16 & โครงการพัฒนาเครื่องสลัดน้ำผักตามอัตลักษณ์ของศูนย์พัฒนาโครงการหลวงทุ่งห & ผศ.เดือนแรม แพ่งเกี่ยว & มทร.ล้านนา พิษณุโลก & 81,500 & 47,780 (59\%)\\
\midrule
\multicolumn{6}{@{}l}{\cellcolor{rpfnavy!12}\bfseries ง8-2 ขับเคลื่อนกลไกการพัฒนาองค์ความรู้เพื่อยกระดับคุณภาพชีวิต (14 โครงการ)}\\
\midrule
1 & โครงการพัฒนาพื้นที่เรียนรู้ศาสตร์พระราชา เพื่อยกระดับศักยภาพพื้นที่ด้ว & นายสามารถ สาลี & กลุ่มแผนงานใต้ร่มพระบารมี & 400,000 & 321,150 (80\%)\\
2 & โครงการการถ่ายทอดองค์ความรู้เรื่องการจัดการขยะแบบยั่งยืน ระดับชุมชนและ & ดร.วนิดา สุริยานนท์ & คณะวิศวกรรมศาสตร์ & 50,000 & 3,090 (6\%)\\
3 & โครงการอบรมเชิงปฏิบัติการถ่ายทอดองค์ความรู้เทคนิคการสร้างแม่พิมพ์เพื่อ & นายคเชนทร์ เครือสาร & สถาบันถ่ายทอดเทคโนโลยีสู่ชุมชน & 50,000 & 47,751 (96\%)\\
4 & โครงการการพัฒนาแหล่งเรียนรู้และถ่ายทอดองค์ความรู้ เทคโนโลยีการพัฒนาผลิ & นายศรีธร อุปคำ & สถาบันถ่ายทอดเทคโนโลยีสู่ชุมชน & 200,000 & 111,400 (56\%)\\
5 & โครงการบริหารจัดการโครงการผลิตเมล็ดพันธุ์และปรับปรุงพันธุ์พืชให้กับศูน & นายศิริชัย แซ่ท้าว & สถาบันวิจัยเทคโนโลยีเกษตร & 100,000 & 45,158 (45\%)\\
6 & โครงการการจัดนิทรรศการเพื่อน้อมถวายรายงานในการรับเสด็จสมเด็จพระกนิษฐาธ & นางสาวสุภัควดี พิมพ์มาศ & มทร.ล้านนา น่าน & 250,000 & 250,000 (100\%)\\
7 & โครงการพัฒนาศักยภาพการเรียนรู้ ด้านการจัดการภูมิทัศน์ ผ่านประสบการณ์จร & นายภาณุพงศ์ สิทธิวุฒิ & มทร.ล้านนา น่าน & 129,050 & 129,050 (100\%)\\
8 & โครงการสนับสนุนการดำเนินงานศูนย์พัฒนาพันธุ์พืชจักรพันธ์เพ็ญศิริ อ.แม่ส & รศ.วิเชษฐ ทิพย์ประเสริฐ & มทร.ล้านนา เชียงราย & 450,000 & 298,151 (66\%)\\
9 & โครงการพัฒนาแหล่งเรียนรู้เกษตรอันเนื่องมาจากพระราชดำริ วิถีธรรมชาติ & รศ.วิเชษฐ ทิพย์ประเสริฐ & มทร.ล้านนา เชียงราย & 118,960 & 15,010 (13\%)\\
10 & โครงการจัดทำคำขอสิ่งบ่งชี้ทางภูมิศาสตร์ไทย (GI) สินค้ากาแฟเลอตอ (อาราบ & ผู้ช่วยศาสตราจารย์ทนงศัก & มทร.ล้านนา ตาก & 36,000 & 0 (0\%)\\
11 & โครงการพัฒนาพื้นที่ต้นแบบ โคก หนอง นา เพื่อการเรียนรู้ด้านเกษตรอินทรีย & ผู้ช่วยศาสตราจารย์ทนงศัก & มทร.ล้านนา ตาก & 65,000 & 0 (0\%)\\
12 & โครงการประยุกต์ใช้การคำนวนสูตรอาหารไก่ไข่ต้นทุนต่ำด้วยโปรแกรมสำเร็จรูป & นางสาวอุษณีย์ภรณ์ สร้อยเ & มทร.ล้านนา พิษณุโลก & 50,000 & 50,000 (100\%)\\
13 & โครงการถ่ายทอดองค์ความรู้และนวัตกรรมการยกระดับวัสดุเพาะเห็ดเศรษฐกิจจาก & นายภควัต หุ่นฉัตร์ & มทร.ล้านนา พิษณุโลก & 50,000 & 40,482 (81\%)\\
14 & โครงการเพิ่มมูลค่าสินค้าเกษตรและการพัฒนาสื่อด้วย AI พื้นที่เขตปฏิรูปที & ผศ.สุริยาพร นิพรรัมย์ & มทร.ล้านนา พิษณุโลก & 50,000 & 49,292 (99\%)\\
\midrule
\multicolumn{6}{@{}l}{\cellcolor{rpfnavy!12}\bfseries ง8-3 พัฒนากำลังคน สร้างอาชีพ ลดความเหลื่อมล้ำ บนพื้นที่สูง (31 โครงการ)}\\
\midrule
1 & โครงการพัฒนาการบริหารจัดการฐานข้อมูลด้วยกลไกการติดตามและประเมินผล โครง & พิมลพรรณ เลิศบัวบาน & กลุ่มแผนงานใต้ร่มพระบารมี & 500,000 & 354,316 (71\%)\\
2 & โครงการพัฒนาฐานข้อมูลด้านวิศวกรรมและยกระดับทักษะอาชีพในพื้นที่โครงการห & นายสามารถ สาลี & กลุ่มแผนงานใต้ร่มพระบารมี & 200,000 & 81,520 (41\%)\\
3 & โครงการยกระดับการบริหารจัดการ ระบบติดตามและเครือข่าย ความร่วมมือกลุ่มแ & นายนริศ กำแพงแก้ว & กลุ่มแผนงานใต้ร่มพระบารมี & 500,000 & 254,112 (51\%)\\
4 & โครงการการจัดกิจกรรมการเรียนรู้ส่งเสริมงานวิจัยและบริการวิชาการ ด้านนว & นายวัชระ กิตติวรเชฏฐ์ & กลุ่มแผนงานใต้ร่มพระบารมี & 60,563 & 60,277 (100\%)\\
5 & โครงการการฝึกอบรมสร้างความสามารถการออกแบบระบบจ่ายน้ำในแปลงเกษตร & นายวัชระ กิตติวรเชฏฐ์ & กลุ่มแผนงานใต้ร่มพระบารมี & 150,000 & 117,038 (78\%)\\
6 & โครงการพัฒนาทักษะอาชีพด้านพลังงานสะอาด & นายวธัญญู วรรณพรหม & กลุ่มแผนงานใต้ร่มพระบารมี & 100,000 & 27,790 (28\%)\\
7 & โครงการนำร่อง SkillChain มทร.ล้านนา :ระบบรับรองทักษะช่างและจัดหางานบน & นางสาวพิมลพรรณ เลิศบัวบา & กลุ่มแผนงานใต้ร่มพระบารมี & 100,000 & 0 (0\%)\\
8 & โครงการการสำรวจผลผลิตทางการเกษตรโครงการหลวง เพื่อการประกอบอาหารและเผยแ & นายวรัญญ์คมน์ รัตนะมงคลช & คณะบริหารธุรกิจและศิลปศาสตร์ & 40,000 & 0 (0\%)\\
9 & โครงการยกระดับมาตรฐานผลิตภัณฑ์ผ้าใยกัญชงวิสาหกิจชุมชนดาวม่างสู่ มผช. ด & ผศ.ญาณิศา โกมลสิริโชค & คณะศิลปกรรมและสถาปัตยกรรมศาสตร์ & 100,000 & 54,144 (54\%)\\
10 & โครงการยกระดับห่วงโซ่คุณค่าห้อมพื้นถิ่น สู่ผลิตภัณฑ์ย้อมสี ธรรมชาติสร้ & ดร.จุฑามาศ ดอนอ่อนเบ้า & คณะศิลปกรรมและสถาปัตยกรรมศาสตร์ & 100,000 & 42,500 (43\%)\\
11 & โครงการการพัฒนานวัตกรรมด้านการศึกษาเพื่อการพัฒนาศักยภาพเยาวชนในโรงเรีย & นายศุภกาล ตุ้ยเต็มวงศ์ & วิทยาลัยเทคโนโลยีและสหวิทยาการ & 149,060 & 89,800 (60\%)\\
12 & โครงการประเมินผลกระทบของโครงการบริการวิชาการมหาวิทยาลัยเทคโนโลยีราชมงค & ผศ.ธัญวดี สุจริตธรรม & สถาบันถ่ายทอดเทคโนโลยีสู่ชุมชน & 142,800 & 120,890 (85\%)\\
13 & โครงการพัฒนาแผนการตลาดของกลุ่มวิสาหกิจชุมชนภายใต้แนวคิด “กินได้ ใช้ได้ & ผศ.ธัญวดี สุจริตธรรม & สถาบันถ่ายทอดเทคโนโลยีสู่ชุมชน & 63,000 & 33,197 (53\%)\\
14 & โครงการต้นแบบอาคารอนุรักษ์พลังงานในรูปแบบ (Green Building \& Green Offi & นายวธัญญู วรรณพรหม & สถาบันถ่ายทอดเทคโนโลยีสู่ชุมชน & 250,000 & 122,500 (49\%)\\
15 & โครงการพัฒนาทักษะวิชาชีพให้กับเจ้าหน้าที่ ชาวบ้าน เยาวชน พื้นที่โดยรอบ & นายศราวุธ วรรณวิจิตร & สถาบันถ่ายทอดเทคโนโลยีสู่ชุมชน & 100,000 & 62,005 (62\%)\\
16 & โครงการพัฒนาต้นแบบชุดน้ำชาเซรามิกสโตนแวร์สำหรับยกระดับผลิตภัณฑ์ชาโครงก & นายสิงหล วิชายะ & สถาบันถ่ายทอดเทคโนโลยีสู่ชุมชน & 100,000 & 50,000 (50\%)\\
17 & โครงการการถ่ายทอดองค์ความรู้กระบวนการผลิตน้ำดื่มสะอาดด้วยระบบรถน้ำดื่ม & นายเมธัส ภัททิยธนี & สถาบันถ่ายทอดเทคโนโลยีสู่ชุมชน & 150,000 & 10,483 (7\%)\\
18 & โครงการพัฒนาทักษะวิชาชีพให้กับเจ้าหน้าที่ ชาวบ้าน เยาวชน พื้นที่โดยรอบ & นายวธัญญู วรรณพรหม & สถาบันถ่ายทอดเทคโนโลยีสู่ชุมชน & 100,000 & 73,386 (73\%)\\
19 & การประยุกต์ใช้ lora mesh ชนิด multi-hops ในการส่งสัญญาณจากเซนเซอร์ตรวจ & นายณัฐวัฒน์ พัลวัล & กลุ่มแผนงานใต้ร่มพระบารมี & 100,000 & 72,800 (73\%)\\
20 & การประยุกต์ใช้เครือข่าย lora-mesh ชนิด multi-hops เพื่อการส่งสัญญาณเตื & นายณัฐวัฒน์ พัลวัล & กลุ่มแผนงานใต้ร่มพระบารมี & 100,000 & 57,500 (57\%)\\
21 & โครงการพัฒนาและยกระดับสมรรถนะด้านการติดตั้งและบำรุงรักษาระบบเซลล์แสงอา & นายวธัญญู วรรณพรหม & กลุ่มแผนงานใต้ร่มพระบารมี & 18,900 & 18,900 (100\%)\\
22 & โครงการวิจัยความเข้มแข็งทางสังคม ชุมชนในการวางแผนการใช้น้ำศูนย์พัฒนาโค & นายวัชระ กิตติวรเชฏฐ์ & กลุ่มแผนงานใต้ร่มพระบารมี & 100,000 & 26,940 (27\%)\\
23 & โครงการส่งเสริมบำรุงรักษาระบบผลิตไฟฟ้าแบบผสมผสานระหว่างระบบผลิตไฟฟ้าพล & นายศรีธร อุปคำ & สถาบันถ่ายทอดเทคโนโลยีสู่ชุมชน & 49,990 & 48,150 (96\%)\\
24 & โครงการพัฒนาทักษะวิชาชีพหลักสูตรการทำมุ้งลวดให้กับสามเณร และเยาวชน พื้ & นายวธัญญู วรรณพรหม & สถาบันถ่ายทอดเทคโนโลยีสู่ชุมชน & 49,860 & 28,860 (58\%)\\
25 & โครงการขับเคลื่อนการสร้างกลไกการส่งเสริมและพัฒนาศักยภาพการเรียนรู้ด้าน & นายรัตนพล พนมวัน ณ อยุธย & สถาบันวิจัยเทคโนโลยีเกษตร & 50,000 & 15,094 (30\%)\\
26 & โครงการพัฒนาทักษะอาชีพและองค์ความรู้ด้านเกษตรและเทคโนโลยีอย่างยั่งยืน & นายก้องเกียรติ ธนะมิตร & มทร.ล้านนา น่าน & 94,000 & 94,000 (100\%)\\
27 & โครงการถ่ายทอดเทคโนโลยีการเพาะขยายพันธุ์ปลาเลียหินเพื่อการอนุรักษ์ปลาท & นายจุลทรรศน์ คีรีแลง & มทร.ล้านนา น่าน & 82,400 & 59,200 (72\%)\\
28 & โครงส่งเสริมพัฒนาโจทย์วิจัยจากชุมชน & รศ.วิเชษฐ ทิพย์ประเสริฐ & มทร.ล้านนา เชียงราย & 100,000 & 8,240 (8\%)\\
29 & โครงการถ่ายทอดองค์ความรู้และนวัตกรรมเพื่อยกระดับมูลค่าไก่ดำ (ไก่ฟ้าหลว & ผศ.ดร.จิรพัฒน์พงษ์ เสนาบ & มทร.ล้านนา เชียงราย & 100,000 & 7,040 (7\%)\\
30 & โครงการพัฒนาระบบติดตามข้อมูลภายใต้สถานีตรวจวัดสภาพอากาศเพื่อป้องกันภัย & ผศ.กิ่งกาญจน์ ปวนสุรินทร & มทร.ล้านนา เชียงราย & 100,000 & 45,000 (45\%)\\
31 & โครงการจัดทำคำขอสิ่งบ่งชี้ทางภูมิศาสตร์ไทย (GI) สินค้ากาแฟเลอตอ (อาราบ & ผู้ช่วยศาสตราจารย์ทนงศัก & มทร.ล้านนา ตาก & 149,000 & 0 (0\%)\\
\midrule
\bottomrule
\end{longtable}
}
